
\subsection{Two-node plumbing and the triple-triangle families}

\begin{result}{B.24}{Proved}{A complete specified realization of six triple-triangle families}
The realization of B.22 extends by the three object families
$\mathrm{III}_2,\mathrm{III}_1,\mathrm{III}'_2$ and the six
arrow families
\[
 \tau_{\pi,\rho},\quad\gamma_{\pi,\rho},\quad
 \tau_r^1,\quad\gamma_r^1,\quad
 \gamma_r^2,\quad\gamma_r^{2'}.
\]
The explicit choices below give $422$ objects and $1020$ arrows.
There are $599$ cycles after a spanning tree and $600$ defining
relators.  The resulting comparison to $\Pi_1M_{0,5}$ is an
equivalence.  Every one of the $660$ added path images and
$420$ added relators is computed in
\texttt{b\_second\_tier.json}.  All arrows of this realization
have explicit expressions through the core paths of B.21 and
the stated twists and one-node erasures.

The triple-triangle twist is a positive quarter-turn.  The
supplied transcription's $1/3$ is not the reading used here:
the enlarged original scan 31 (folio 30), independently viewed
by two readers, supports $1/4$ and matches the upper quarter-turn
glyph.  The alternative transcription is retained in the reading
register.  The endpoint and path conventions below are part of
the realization; no additional unrecorded direct-path convention
is inferred from an endpoint count.
\end{result}
\begin{proof}
Put $\epsilon=1/10$ and label the following five-point
configurations by $0,1,2,3,4$ in the displayed order:
\[
\begin{aligned}
 R&=(\infty,-1-\epsilon,1+\epsilon,-1+\epsilon,1-\epsilon),\\
 M&=(\infty,-1-i\epsilon,1+\epsilon,-1+i\epsilon,1-\epsilon),\\
 L&=(\infty,-1-i\epsilon,1+i\epsilon,-1+i\epsilon,1-i\epsilon).
\end{aligned}                                                     \tag{B.54}
\]
The two node clusters are $A=\{1,3\}$ and $B=\{2,4\}$;
the remaining mark is $0$.  Near their centres $-1,1$ use
\[
 z_A=z+1,\quad z_B=z-1,\qquad
 z_Av_A=q_A,\quad z_Bv_B=q_B.
\]
The two three-point tails have coordinates with marks
$-1,1,\infty$.  Thus (B.54) are respectively the smoothing
parameters $(\epsilon,\epsilon)$,
$(i\epsilon,\epsilon)$ and $(i\epsilon,i\epsilon)$.
Discs of radius $1/4$ about $-1,1$ are disjoint and contain
the corresponding pairs.  Choose the two real rays on each
untwisted seam; after a quarter-turn retain their union with
the rotated two rays.  This gives the $2+2$, $4+2$ and
$4+4$ seam markings of the three diagrams.

The holomorphic stabilizers of $R,M,L$ are, respectively,
\[
 H=\langle(12)(34)\rangle,\qquad 1,\qquad H.
                                                               \tag{B.55}
\]
For $R,L$ the nonidentity element is $z\mapsto-z$.
To exclude any additional automorphism, a permutation of
the five labels determines at most one fractional linear map.
Normalize its first, second and fifth labels to $0,1,\infty$
and compare the remaining two cross-ratios.  The exact rational
calculation for all $120$ permutations gives exactly (B.55);
the finite test and all configurations are in the program.
Hence the object sets are
$\mathfrak S_5/H,\mathfrak S_5,\mathfrak S_5/H$, of sizes
$60,120,60$.

Here are all six representative paths.  In the first two,
replace the positive circular quarter-turn by its chord
\[
 u(t)=\epsilon(1-t+it),\qquad 0\leq t\leq1.
\]
The radial interpolation to the circular path avoids zero
inside the relevant disc and fixes all other marks.  The
two twist paths are
\[
\begin{aligned}
 T_A(t)&=(\infty,-1-u(t),1+\epsilon,-1+u(t),1-\epsilon),\\
 T_B(t)&=(\infty,-1-i\epsilon,1+u(t),-1+i\epsilon,1-u(t)).
\end{aligned}                                                     \tag{B.56}
\]
Their endpoints are $R\to M$ and $M\to L$.

Define the two target configurations
\[
\begin{aligned}
 R_*&=(\infty,-(1-\epsilon)^{-1},0,
                 -(1+\epsilon)^{-1},-1/2),\\
 M_*&=(\infty,-(1-i\epsilon)^{-1},0,
                 -(1+i\epsilon)^{-1},-1/2).
\end{aligned}                                                     \tag{B.57}
\]
The path $G_R$ is the straight coordinate interpolation
from $R$ to $R_*$, and $G_M$ is that from $M$ to $M_*$,
keeping the first coordinate $\infty$.  In $G_R$ the real
ordering $1<3<4<2$ is preserved.  In $G_M$ the conjugate
pair $\{1,3\}$ stays off the real line and the two other
finite marks remain real and ordered.  Thus both avoid
collisions.  More precisely the retained $A$ pair has
half-separation at most $10/99$; its centre is separated
from the other finite marks by more than $1/2$.  A moving
disc of radius $1/4$ therefore retains that pair throughout.
The complement has four marked boundary or puncture positions.
Consequently these paths are actual one-node erasures with
a fixed three-point tail, including its displayed framing.
At the endpoints use the change of coordinate $z\mapsto
-1/(z+1)$ to compare the retained disc with the one-node
model of B.22.  Its restriction transports the cut and all
its pins; a smaller disc and the intervening annulus give
the same isotopy class.

For the last two paths put $s=1-t$ and
$v=\epsilon+(1-\epsilon)t$.  Set
\[
\begin{aligned}
 G_L(t)&=(\infty,-s-i\epsilon,s+iv,-s+i\epsilon,s-iv),\\
 G_{LL}(t)&=(\infty,-1-iv,1+iv,-1+iv,1-iv).
\end{aligned}                                                     \tag{B.58}
\]
For $t<1$ in $G_L$ the two pair centres have different
real parts; at $t=1$ the imaginary magnitudes are
$\epsilon$ and $1$, so there is still no collision.
The pair $A$ remains in a disc of radius $1/4$ around
$-s$: indeed the distance from this centre to either
other point is
\[
 \sqrt{4(1-t)^2+(\epsilon+(1-\epsilon)t)^2}>1/2.
\]
Thus $G_L$ also carries a specified retained cut.  When
comparing its endpoint with the two-pin right object, use
the two pins selected by the transported normal ray;
the extra pins are explicitly erased.  This is the combined
circle and marking erasure convention for this realization.
The path $G_{LL}$ erases both cuts and reaches a square.
It fixes the real parts $-1,1$, so again no collision occurs.

To state every endpoint without an implicit relabelling,
permutations act on the labels on the left.  Write
\[
\begin{aligned}
 p_R&=(2,1,3,0,4),&
 p_L&=(0,1,3,4,2),&
 p_Q&=(0,4,3,1,2),
\end{aligned}                                                     \tag{B.59}
\]
where tuples list the images of $0,1,2,3,4$.
Then (B.57) are $p_R V(\epsilon)$ and
$p_R V(i\epsilon)$ after $z\mapsto-1/(z+1)$.
At the end of $G_L$, multiplication by $-i$ identifies
the configuration with $p_L V(\epsilon)$.  At the end
of $G_{LL}$, multiplication by $(1-i)/2$ identifies it
with $p_QK(1)$.  This proves the following endpoint table:
\[
\begin{array}{c|c|c|c|r}
 \text{family}&\text{parameter}&\text{source}&\text{target}&
 \text{arrows}\\ \hline
 \tau_{\pi,\rho}&g\in\mathfrak S_5&gH\ (R)&g\ (M)&120\\
 \gamma_{\pi,\rho}&g\in\mathfrak S_5&gH\ (R)&gp_RH\ (\mathrm{II}_2)&120\\
 \tau_r^1&g\in\mathfrak S_5&g\ (M)&gH\ (L)&120\\
 \gamma_r^1&g\in\mathfrak S_5&g\ (M)&gp_RH\ (\mathrm{II}'_2)&120\\
 \gamma_r^2&g\in\mathfrak S_5&gH\ (L)&gp_LH\ (\mathrm{II}_2)&120\\
 \gamma_r^{2'}&gH\in\mathfrak S_5/H&gH\ (L)&gp_QH_Q\ (\mathrm I_4)&60.
\end{array}                                                     \tag{B.60}
\]
In the last row $p_Q^{-1}Hp_Q\subset H_Q$, and the
path is equivariant for $H$, so its coset parameter is
well defined.  The other five rows have trivial arrow
stabilizer.  The table proves all endpoint valencies
directly and records the three permutations that would
otherwise be hidden in a comparison of diagrams.

Keep B.22's indices, then number the new objects $182$--$241$,
$242$--$361$ and $362$--$421$ in lexicographic coset or
permutation order.  Append the arrows in the order (B.60),
also lexicographically.  Normalize each of their five-point
paths by (B.41) and use the same base configuration and
least-integer detour transports as B.21.  The computation
certifies $724$ rational subintervals by exact Sturm tests.
Every pair difference stays in a strict half-plane on each
subinterval; straight interpolation therefore supplies an
actual collision-free homotopy to the polygon used for
the braid calculation.  Crossing signs use the same exact
projection of slope $1/37$.  All $660$ resulting pure-braid
words are combed and checked by the faithful Artin action.

Extend the old spanning tree greedily by these arrows.
It now has $421$ edges.  It is already connected before
the last three rows of (B.60); its new edges are twists
and $G_R$ erasures.  The graph has
$1020-422+1=599$ free cycles.  For each of its $420$ new
non-tree cycles impose its computed word in the eleven
core cycles, using B.21's inverse dictionary.  Together
with the old $180$ relators this gives $600$ relators.
All $420$ expanded equations are retained and independently
checked in the output.  Eliminating their new cycle
generators reduces the presentation exactly to B.22.
The actual path functor, its inverse on the eleven core
cycles, and the tree transports therefore prove the
claimed equivalence.

Finally, the same equations express every added arrow
as a tree transport, a word in the core cycles, and an
inverse tree transport.  The new tree transports are
the twists and the displayed one-node erasures.  The
core paths have the exact four-point covering construction
of B.23.  This gives the stated expressions through
these specified relative-dimension-one operations,
including the simultaneous two-cut erasure $G_{LL}$.
The remaining tetrahedron--triangle families, and any
further comparison with unstated choices in the manuscript,
are not supplied by this argument.
\end{proof}
